\chapter{Resumen Temario}
% TEMA 1: La economía española en perspectiva
\section*{Tema 1: La economía española en perspectiva: de la autarquía a la convergencia europea}

\subsection*{1. Crecimiento y convergencia de la economía española (1959-2024)}

\textbf{El punto de partida: la autarquía (1939-1959)}

\begin{itemize}
  \item Este periodo se caracterizó por un fuerte intervencionismo gubernamental y un severo aislamiento exterior.
  \item Esta situación generó desequilibrios insostenibles: déficit en la balanza por cuenta corriente, espiral de precios-salarios desde 1955 y déficit público.
  \item A nivel internacional, contrastaba con el nuevo orden geopolítico posterior a la Segunda Guerra Mundial, la creación de organismos como el FMI, el Banco Mundial y el GATT, y la consolidación de la economía de mercado en Europa.
\end{itemize}

\textbf{El Plan de Estabilización de 1959}

Ante la quiebra inminente, se aprobó este plan con resultados muy positivos, estructurado en tres pilares:

\begin{enumerate}
  \item \textbf{Medidas de apoyo:} Consenso interno, gobierno sólido (1957) y apoyo financiero de organismos internacionales (FMI, OECE).
  \item \textbf{Políticas de ajuste:} Búsqueda del equilibrio presupuestario y monetario interior (reforma fiscal, contención del gasto y de la demanda) y exterior (integración de la peseta en el Sistema Monetario Internacional, apertura a la inversión extranjera y fomento de exportaciones).
  \item \textbf{Reformas estructurales:} Adhesión a organismos internacionales, nuevo arancel (1960), racionalización administrativa y una tímida flexibilización del sistema productivo y laboral (seguro de desempleo en 1961).
\end{enumerate}

\textbf{Rasgos generales del crecimiento (1961-2024)}

\begin{itemize}
  \item España ha crecido a un ritmo superior al europeo (2,4\% de media anual acumulativa en PIB per cápita frente al 2,0\% de la Eurozona).
  \item El perfil temporal de crecimiento ha sido sincrónico al de los países occidentales.
  \item Sin embargo, España sufre las crisis económicas con mayor intensidad y presenta fluctuaciones del PIB más marcadas, debido a una liberalización tardía y rigideces estructurales persistentes.
\end{itemize}

Se divide este crecimiento en cinco grandes etapas:

\subsection*{2 y 3. Las fases de crecimiento, convergencia y crisis económicas}

\textbf{1.\textsuperscript{a} Etapa (1960-1974): El Desarrollo}

\begin{itemize}
  \item Se implementaron los ``Planes de Desarrollo'' (cuatrienales, vinculantes para el sector público e indicativos para el privado).
  \item Favorecido por la expansión internacional, generó un gran crecimiento del PIB, progresiva desagrarización y apertura exterior.
  \item \textit{Desequilibrios:} crecimiento desequilibrado, inflación media del 7\% y déficit exterior.
\end{itemize}

\textbf{2.\textsuperscript{a} Etapa (1974-1985): Crisis económica y transición}

\begin{itemize}
  \item Marcada por la crisis internacional del petróleo (1974) y la transición política, provocando divergencia en renta con Europa.
  \item Generó caída productiva, aumento del paro, espiral precios-salarios y desequilibrio externo.
  \item \textbf{Pactos de la Moncloa (1977):} programa de saneamiento y reforma consensuado políticamente para reducir la inflación y reestructurar la economía.
  \item Lograron reducir desequilibrios temporalmente, pero a partir de 1978 la inestabilidad política (dimisión de Suárez, intento de golpe de Estado de 1981) paralizó las reformas.
\end{itemize}

\textbf{3.\textsuperscript{a} Etapa (1985-1995): Integración Europea y crisis de los 90}

\begin{itemize}
  \item El gobierno socialista (1982) priorizó crear empleo frente a graves problemas estructurales (crisis industrial, rigidez laboral, déficit de empresas públicas y Seguridad Social).
  \item \textbf{Adhesión a la CEE (1 enero 1986):} inició un ciclo de convergencia con un crecimiento del PIB del 4,4\% (1986-1991), impulsado por bajos tipos, entrada de capital extranjero y un contexto internacional favorable (caída del petróleo y dólar).
  \item \textit{Gestación de desequilibrios:} inflación dual y diferencial frente a Europa (7\%), déficit público estructural, sobrevaloración de la peseta (en la parte alta de la banda del Sistema Monetario Europeo) y déficit comercial.
  \item \textbf{Tratado de Maastricht (1992):} definió las fases hacia la Unión Económica y Monetaria (UEM) y los criterios de convergencia nominal (límites estrictos en inflación, tipos de interés, déficit < 3\% y deuda < 60\%).
  \item \textbf{Crisis de 1992-1994:} el colapso del SME (altos tipos alemanes por la reunificación), la sobrevaloración de la peseta y la falta de competitividad forzaron tres devaluaciones de la moneda (20\%). Produjo destrucción de empleo y déficits públicos superiores al 7\%.
  \item Mediante los \textit{Programas de Convergencia} (especialmente el de 1994/1997), la Ley de Autonomía del Banco de España (1994) y privatizaciones, España logró enderezar el rumbo y cumplir los criterios.
\end{itemize}

\textbf{4.\textsuperscript{a} Etapa (1996-2007): Expansión sostenida}

\begin{itemize}
  \item Periodo de gran estabilidad macroeconómica regido por los \textit{Programas de Estabilidad} (Pacto de Estabilidad y Crecimiento de la UE) donde la política monetaria pasó al BCE.
  \item \textit{Factores expansivos:} creación de 7,84 millones de empleos (tasa de paro bajó al 8,3\%), auge del consumo (efecto riqueza, bajos tipos) e inmensa inversión en capital fijo impulsando una \textbf{burbuja inmobiliaria}.
  \item Se logró superávit público (2005-2007) y reducir la deuda al 36\% del PIB.
  \item \textit{Problemas ocultos bajo la bonanza:} deterioro brutal del déficit exterior (10\% del PIB en 2007) por la inflación diferencial y el estancamiento de la productividad. La economía se especializó en sectores de tecnología media-baja, el mercado laboral se dualizó, el sistema financiero asumió riesgos extremos ligados al ``ladrillo'' y el sector público ejecutó infraestructuras innecesarias sin reformar las pensiones.
\end{itemize}

\textbf{5.\textsuperscript{a} Etapa (2008-2024): Gran Recesión, COVID-19 y recuperación}

\begin{itemize}
  \item \textbf{La crisis financiera (2008-2014):} detonada externamente (crisis \textit{subprime} en EE. UU., petróleo) e internamente por el estallido de la burbuja inmobiliaria y los desequilibrios acumulados (exceso de liquidez, dependencia exterior, endeudamiento a corto plazo para financiar el largo plazo).
  \item \textit{Impacto:} el crédito se contrajo. El déficit público se disparó al 11,3\% (2009) y la deuda superó el 100\% (2014). El PIB real cayó un 8,6\% y la tasa de paro alcanzó el 26,1\% en 2013 (destrucción masiva por la rigidez institucional y contratación temporal).
  \item \textbf{Recuperación (2014-2019):} crecimientos del PIB superiores al 3\%, apoyados en exportaciones, superávit por cuenta corriente, desapalancamiento privado y saneamiento bancario. Sin embargo, persistía la precariedad laboral (27\% de temporalidad) y la alta deuda pública.
  \item \textbf{Crisis sanitaria COVID-19 (2020):}
    \begin{itemize}
      \item \textit{Causas del grave impacto en España:} reacción tardía siendo potencia turística, alta densidad y transporte público en ciudades clave (Madrid, Barcelona), cohabitación intergeneracional, envejecimiento y severas carencias estructurales sanitarias (solo 3 camas/100.000 hab. frente a las 14 de Japón, falta de reservas estratégicas).
      \item \textit{Efectos económicos:} doble \textit{shock} de oferta y demanda global. España fue extremadamente vulnerable por su dependencia del turismo (9,2\% del PIB) y la automoción, alta proporción de microempresas/autónomos, y su temporalidad laboral.
      \item \textit{Respuesta nacional:} avales ICO, aplazamiento de impuestos, agilización de ERTEs por fuerza mayor y creación del Ingreso Mínimo Vital (IMV).
      \item \textit{Respuesta europea:} a diferencia de 2008, la UE aplicó máxima expansión: liquidez del BEI, fondos MEDE, programa PEPP del BCE (1,3 billones) e inyección directa mediante el plan \textit{Next Generation EU} (750.000 millones de euros).
    \end{itemize}
\end{itemize}

\subsection*{4. Balance del periodo y retos pendientes}

Si bien la transformación económica y social desde 1960 ha sido extraordinaria, el documento subraya que España debe afrontar con urgencia reformas inaplazables:

\begin{itemize}
  \item Sanear y equipar el sistema sanitario.
  \item Articular un ecosistema real de innovación ($I+D+i$).
  \item Potenciar la educación y la FP Dual.
  \item Reformar el sistema de pensiones para garantizar su viabilidad ante el reto demográfico.
  \item Erradicar la dualidad y precariedad laboral.
  \item Abordar decididamente la transición ecológica y la digitalización.
\end{itemize}

% Le recomiendo integrar la interrelación entre los ciclos económicos descritos y los desequilibrios estructurales crónicos (déficit público, dualidad laboral y baja productividad) de cara a su estudio, pues constituyen la columna vertebral de la asignatura.

\clearpage
% TEMA 2: Factores productivos y fuentes de crecimiento
\section*{Tema 2: Factores productivos y fuentes de crecimiento en la economía española}

\subsection*{1. Transformaciones estructurales}

El desarrollo de la economía española a largo plazo viene indisolublemente ligado a un profundo cambio en la estructura sectorial del empleo.

\begin{itemize}
  \item \textbf{Desagrarización y tercerización:} En 1960, el sector primario (agricultura y pesca) aglutinaba el 38,7\% del empleo nacional, cifra que se ha desplomado hasta un marginal 3,6\% en 2023.
  \item Simultáneamente, la industria, energía y construcción pasaron de representar el 30,3\% en 1960 al 18,2\% en la actualidad.
  \item Como contrapartida, el sector servicios ha experimentado un crecimiento colosal, pasando de emplear al 31,0\% de la población ocupada en 1960 a un dominante 78,2\% en 2023. Esta transición alinea la estructura española con la de las economías más avanzadas de nuestro entorno.
\end{itemize}

\subsection*{2. Determinantes del crecimiento a largo plazo}

El crecimiento económico a largo plazo se fundamenta en el aumento de la renta per cápita, el cual es el resultado multiplicativo de dos grandes determinantes: la productividad del trabajo y el empleo per cápita.

\[
\frac{\text{PIB}}{\text{Población}} = \left(\frac{\text{PIB}}{\text{Empleo}}\right) \times \left(\frac{\text{Empleo}}{\text{Población}}\right)
\]

\begin{itemize}
  \item \textbf{Productividad del trabajo ($PIB/Empleo$):} Mide el rendimiento medio de los trabajadores.
  \item \textbf{Empleo per cápita ($\text{Empleo}/\text{Población}$):} Refleja el porcentaje de la población total que realiza actividades productivas.
  \item La tasa de variación de la renta per cápita se aproxima sumando las tasas de variación de ambas ratios.
\end{itemize}

Al analizar el periodo 1961-2024, el PIB per cápita en España creció a una tasa media anual del 2,4\%, sustentado fundamentalmente en los avances de la productividad del trabajo (2,3\%), mientras que el empleo per cápita apenas aportó un 0,1\% en el cómputo global de dicho periodo.

Para desglosar el empleo per cápita, debemos observar la demografía:

\[
\frac{\text{Empleo}}{\text{Población}} = \left(\frac{\text{Empleo}}{\text{PET}}\right) \times \left(\frac{\text{PET}}{\text{Población}}\right)
\]

\textit{(Donde PET = Población en Edad de Trabajar)}.

El crecimiento de la PET fue determinante. A diferencia de Estados Unidos y Europa Occidental donde despuntó en los años 60, en España (y también en Irlanda) este ensanchamiento poblacional fue más tardío y prolongado. Esto se explica por el retraso de nuestro \textit{baby boom} y, posteriormente, por una intensa entrada de flujos migratorios. Como evidencia, la PET española suponía el 64,5\% del total en 1960, escaló al 68,9\% en 2005 y se situó en el 65,4\% en 2022.

\subsection*{3. Productividad, capital y progreso tecnológico}

Para comprender la productividad, es imperativo diseccionar los factores que la impulsan.

\textbf{A. Capital Físico}

La Contabilidad Nacional clasifica el capital físico en tres grandes bloques:

\begin{enumerate}
  \item \textbf{Capital Directamente Productivo (CDP):} Generado por el sector privado e incluye maquinaria, equipos de transporte, edificaciones e instalaciones.
  \item \textbf{Infraestructuras y equipamientos:} Proporcionado por el sector público, presta servicios vitales al CDP (redes de transporte, comunicaciones, energía, instalaciones sanitarias y educativas).
  \item \textbf{Capital Residencial:} Comprende las viviendas y no se destina directamente a la producción de bienes y servicios.
\end{enumerate}

\textbf{B. Capital Humano}

Se define como la combinación de salud, conocimientos y destrezas de los trabajadores, traduciéndose en mayor rendimiento laboral y salarial. Es una inversión y un factor estratégico que potencia el factor trabajo, incide en la generación/difusión tecnológica y eleva la competitividad.

\begin{itemize}
  \item Se genera mediante la educación formal, la experiencia/aprendizaje en el puesto, o ambas.
  \item El esfuerzo inversor (Gasto en educación/PIB) oscila entre el 4\% y el 5\% desde los años noventa.
  \item Los resultados muestran una población mucho más cualificada: el número medio de años de estudio saltó de 5,3 años en 1960 a 12,1 años en 2024. No obstante, este enorme salto formativo no se ha correspondido proporcionalmente con la evolución de la Productividad Aparente del Trabajo (PAT).
\end{itemize}

\textbf{C. La Función de Producción y el Residuo de Solow}

El incremento histórico de la PAT en España (1960-2024), cuantificado en +47.000€, se explica por:

\begin{itemize}
  \item Capital físico por trabajador: +17.000€ (36\%).
  \item Capital humano: +7.000€ (15\%).
  \item Productividad Total de los Factores (PTF): +24.000€ (49\%).
\end{itemize}

Analíticamente, esto se explica asumiendo una función Cobb-Douglas:

\[
Y = A K^{\alpha} h^{\beta} L^{(1-\alpha)}
\]

\textit{(Donde $Y$ es producto, $A$ es la PTF o progreso técnico, $K$ es capital físico, $h$ capital humano, $L$ trabajo no cualificado, y $\alpha, \beta$ son las elasticidades)}.

Dividiendo por $L$ para obtener la productividad media del trabajo ($y$):

\[
y = A k^{\alpha} h^{\beta}
\]

Expresado en tasas de variación para analizar el crecimiento:

\[
\Delta y = \Delta A + \alpha\Delta k + \beta\Delta h
\]

Despejando el término $\Delta A$ obtenemos la Productividad Total de los Factores o ``Residuo de Solow'':

\[
\Delta A = \Delta y - \alpha\Delta k - \beta\Delta h
\]

La trayectoria de la PTF en España es muy reveladora: creció de forma sostenida al 2\% anual entre 1975 y 2000, pero desde el inicio del decenio de 1990 su aportación ha sido prácticamente nula, evidenciando un grave estancamiento en la eficiencia combinada de capital y trabajo. La PTF relativa respecto a la Eurozona cayó sostenidamente hasta la Gran Recesión, se desplomó puntualmente en 2020 (debido a la retención artificial de empleo/capital durante la pandemia) y se recupera apreciablemente desde 2021.

\subsection*{4. Fluctuaciones cíclicas y desequilibrios macroeconómicos}

Finalmente, para analizar los desequilibrios coyunturales, debemos distinguir dos variables fundamentales:

\begin{itemize}
  \item \textbf{PIB Potencial:} Describe la tendencia de la oferta agregada a largo plazo, impulsada por demografía y productividad.
  \item \textbf{Output gap (brecha de producción):} Diferencia entre el PIB potencial y el real motivada por \textit{shocks} transitorios de oferta y demanda a corto/medio plazo.
\end{itemize}

En la historia económica española:

\begin{itemize}
  \item \textbf{Fases expansivas (PIB real por encima del potencial):} Abarcan los periodos 1960-1974, 1986-1991, 1995-2007 y el momento de recuperación actual.
  \item \textbf{Fases recesivas (PIB real por debajo del potencial):} Coinciden con la crisis de los años 70, la crisis de 1993, la crisis financiera global de 2009 y el profundo abismo generado por la pandemia de la Covid-19 en 2020.
\end{itemize}
 
\clearpage
% TEMA 3: El sector primario
\section*{Tema 3: El sector primario}

\subsection*{1. Conceptos básicos y rasgos estructurales del sector primario}

El sector primario se define en la Contabilidad Nacional bajo la rama de ``Agricultura, ganadería, selvicultura y pesca''.

\begin{itemize}
  \item \textbf{Composición interna (2023):} La inmensa mayoría del sector corresponde a la agricultura y ganadería, que acaparan el 93,3\% del Valor Añadido Bruto (VAB) y el 92,6\% del empleo. La pesca y acuicultura representan el 3,8\% del VAB y el 5,2\% del empleo, mientras que la silvicultura y explotación forestal suponen un 2,9\% del VAB y un 2,2\% del empleo.
  \item \textbf{El sistema agroindustrial (2022):} La importancia del sector se magnifica al considerar las industrias manufactureras a las que provee materias primas. En este complejo conjunto, el sector primario estricto aporta el 46,6\% del VAB y el 56,5\% del empleo (710.000 puestos). Por su parte, el sector industrial asociado (alimentos, bebidas, transformados de pesca, madera y papel) aporta el 53,4\% del VAB y el 43,5\% del empleo.
  \item \textbf{Orientaciones productivas (2024):} Destacan los granívoros (porcino y avícola) con el 25,2\% de la producción total (usando solo el 2,6\% de la superficie). Le siguen la horticultura (19,3\% de producción en 1,2\% de superficie), los bovinos (14,3\% de producción en 14,5\% de superficie) y los frutales/cítricos (13,2\% de producción en 6,4\% de superficie). El olivar representa el 10,6\% de la producción ocupando el 9,9\% de la tierra.
\end{itemize}

\subsection*{2. Evolución de las macromagnitudes agrarias}

A lo largo de las últimas décadas, el sector agrario ha experimentado una notable pérdida de peso relativo en el conjunto de la economía española, un proceso propio del desarrollo económico.

\begin{itemize}
  \item \textbf{Valor Añadido Bruto (VAB):} En 1985, el VAB agrario suponía el 5,7\% del VAB total a precios corrientes, cayendo paulatinamente hasta situarse en un 2,8\% en 2024. Medido a precios constantes, la caída ha sido del 3,2\% al 2,5\% en el mismo periodo.
  \item \textbf{Empleo:} La pérdida de peso es aún más drástica en términos de ocupación, pasando de representar el 14,1\% del empleo total nacional en 1985 a tan solo un 3,5\% en 2024.
\end{itemize}

\subsection*{3. Comportamiento del sector: factores de demanda y de oferta}

El desempeño del sector primario se explica a través de cuatro grandes vectores analíticos:

\subsubsection*{3.1. Intercambios comerciales con el exterior}

\begin{itemize}
  \item Se observa una creciente integración internacional: el coeficiente de apertura externa se ha disparado desde el 59,7\% en 1985 hasta el 168,1\% en 2024.
  \item España mantiene un sólido superávit comercial agroalimentario. La tasa de cobertura ha pasado del 104,2\% en 1985 al 147,2\% en 2024.
  \item Somos altamente competitivos exportando hortalizas, frutas, bebidas, aceite de oliva, carnes y ganado en pie. Por el contrario, somos deficitarios en tabaco, madera, cereales, semillas oleaginosas y lácteos.
  \item Por mercados, la tasa de cobertura en transacciones intracomunitarias alcanza el 175\%, mientras que en el comercio extracomunitario es deficitaria (93\%), influida por la salida del Reino Unido.
\end{itemize}

\subsubsection*{3.2. Especialización productiva}

\begin{itemize}
  \item El sistema agroindustrial actual es un modelo capitalizado y modernizado, con uso intensivo de consumos intermedios industriales y fuertemente condicionado por las subvenciones de la Política Agrícola Común (PAC).
  \item En el marco de la Unión Europea (2024), España lidera en la producción de olivar (aportando el 59,0\% del total de la UE-27), frutales y cítricos (29,3\%), granívoros (21,3\%) y horticultura (19,8\%). Presentamos una baja especialización en agricultura general (7,9\%) y bovinos (9,4\%).
\end{itemize}

\subsubsection*{3.3. Eficiencia productiva: análisis de la productividad}

La productividad aparente del trabajo agrario se desagrega matemáticamente en dos factores:

\[
\frac{Q}{L} = \frac{Q}{T} \cdot \frac{T}{L}
\]

Donde la productividad de la tierra ($Q/T$, medida como VAB/Hectárea) se multiplica por la ratio de estructuras o superficie por ocupado ($T/L$).

\begin{itemize}
  \item La productividad del trabajo se ha cuadriplicado en términos reales, pasando de 12.993 euros/empleo en 1985 a 52.929 euros/empleo en 2024.
  \item Esto se ha logrado gracias al aumento de la productividad de la tierra (de 878 a 1.802 euros/ha) y, sobre todo, al incremento de la superficie por empleo (de 14,0 a 29,4 hectáreas/empleo).
  \item \textbf{Causas del aumento de productividad:} (1) el fuerte encarecimiento de los salarios agrarios en comparación con la maquinaria, lo que ha incentivado la mecanización y la reducción de mano de obra; y (2) el hecho de que los precios percibidos por los agricultores han crecido menos que el IPC (coste de vida), obligándolos a producir más volumen físico para mantener su renta real.
\end{itemize}

\subsubsection*{3.4. Dimensión económica de las explotaciones y eficiencia}

\begin{itemize}
  \item Existe una clara relación entre el tamaño de la explotación (medido en Unidades de Dimensión Europea - UDE) y su eficiencia.
  \item Las explotaciones pequeñas (<16 UDE) representan el 61,5\% del total, pero apenas aportan el 15,5\% de la Producción Estándar Total (PET).
  \item En el extremo opuesto, las grandes explotaciones (>60 UDE) suponen solo el 17,7\% del total de fincas, pero gestionan el 62,2\% de la superficie y generan un abrumador 80,2\% de la PET. Además, su productividad del trabajo (123.000 euros/UTA) barre estadísticamente a las explotaciones más pequeñas.
\end{itemize}

\subsection*{4. La Política Agrícola Común (PAC) y su reforma}

La PAC nació en 1962 con los objetivos históricos de garantizar la renta agrícola, estabilizar mercados y asegurar el abastecimiento alimentario a precios razonables.

\begin{itemize}
  \item \textbf{Principios orientadores:} Unidad de Mercado (libre circulación), Preferencia Comunitaria (protección exterior) y Solidaridad Financiera (costes compartidos). El instrumento clásico fueron las Organizaciones Comunes de Mercado (OCM), fusionadas en una OCM única en 2007.
  \item \textbf{Reformas recientes:} Tras las reformas de 2003, 2009 y 2013, la actual \textbf{PAC Post-2020 (aplicable 2023-2027)} asume un fuerte enfoque medioambiental, de apoyo a pequeñas explotaciones y mayor flexibilidad estatal.
  \item \textbf{Pacto Verde Europeo (PVE) y estrategia ``De la granja a la mesa'':} Impone una ``ecologización de la PAC'', exigiendo que el 40\% del presupuesto se destine al cambio climático. Fija para 2030 metas estrictas: reducir el uso de fertilizantes y pesticidas, y lograr que la agricultura ecológica ocupe el 25\% de la Superficie Agrícola Utilizada (SAU).
\end{itemize}

\textbf{La nueva ``arquitectura verde'' de la PAC engloba:}

\begin{itemize}
  \item \textbf{Pago básico:} Ahora sujeto a una \textit{condicionalidad reforzada} (requisitos climáticos/ambientales más estrictos).
  \item \textbf{Eco-Esquemas:} Pagos adicionales voluntarios para agricultores que asuman compromisos ambientales superiores.
  \item \textbf{Techo de ayudas (capping):} Modulación de ayudas a partir de 60.000 euros y fijación de un tope máximo de 100.000 euros por explotación (descontando costes laborales).
  \item \textbf{Otras medidas clave:} Pagos específicos para favorecer el relevo generacional y la mujer rural, una reserva de crisis de 1.500 millones de euros, y el mantenimiento de ayudas asociadas para sectores frágiles (ej. ganado herbívoro).
  \item Todo esto se vertebra a nivel estatal mediante la exigencia de un \textbf{Plan Estratégico Nacional}, el cual otorga mayor poder de decisión a los Estados pero requiere validación de la Comisión Europea.
\end{itemize}
 
\clearpage
% TEMA 4: El sector industrial
\section*{Tema 4: El sector industrial}

\subsection*{A. Sector Energético}

El sector energético posee un carácter estratégico transversal, ya que la energía es un \textit{input} fundamental para todas las actividades productivas, especialmente la industria. Su eficiencia genera un fuerte efecto de arrastre sobre el resto de la economía.

\textbf{Rasgos estructurales y evolución:}

\begin{itemize}
  \item \textbf{Dependencia y autoabastecimiento:} Es una rama con vocación importadora (petróleo, gas natural, carbón). El grado de autoabastecimiento nacional es alarmantemente bajo, cubriendo apenas el 30\% de las necesidades de energía primaria. El petróleo sigue siendo dominante (~45,8\% en 2023), aunque ha bajado desde los años setenta.
  \item \textbf{Transformación de la demanda y oferta:} Históricamente se observan procesos de sustitución: del carbón al petróleo (60s-70s), hacia la nuclear (80s), hacia el gas natural (90s) y, actualmente, hacia las energías renovables. En 2023, las renovables ya cubrieron el 50,3\% de la producción eléctrica nacional, desplazando a los ciclos combinados de gas.
  \item \textbf{Características del sector:} Se trata de una actividad altamente intensiva en capital, con una producción nacional insuficiente y una altísima dependencia externa. Además, existe una fuerte internacionalización de las grandes firmas españolas (ej. Repsol en Latinoamérica).
  \item \textbf{Concentración de mercado (oligopolios):} Pese a la privatización de empresas como Repsol o Endesa, el sector mantiene barreras de entrada y alta concentración.
  \item \textit{Hidrocarburos:} Repsol, Cepsa y BP controlan más del 70\% del mercado.
  \item \textit{Gas natural:} Naturgy es el operador dominante, seguido por Iberdrola y Endesa.
  \item \textit{Electricidad:} Iberdrola y Endesa generan cerca del 75\% de la electricidad vendida y concentran la mitad de la potencia instalada. Red Eléctrica de España mantiene el monopolio natural del transporte en alta tensión.
\end{itemize}

\subsection*{B. El sector industrial manufacturero}

\subsubsection*{1. Delimitación y clasificación}

La industria transformadora (manufacturas) presenta una productividad superior a la agricultura y los servicios, actuando como motor de innovación y tracción económica. Las actividades se clasifican cruzando dos variables: el dinamismo de la demanda y la intensidad tecnológica.

\begin{itemize}
  \item \textbf{Avanzadas (alta tecnología y demanda):} Informática, electrónica, óptica, farmacia. Tienen alto esfuerzo innovador, personal muy cualificado y mayor capital extranjero.
  \item \textbf{Intermedias (media-alta tecnología y demanda):} Química, maquinaria, material eléctrico y vehículos de motor.
  \item \textbf{Tradicionales (baja/media-baja tecnología y demanda):} Alimentación y bebidas, textil, calzado, madera, metalurgia. Operan con demanda débil, menor esfuerzo innovador y mayor uso de trabajo por unidad de producto.
\end{itemize}

\subsubsection*{2. Evolución del sector}

\begin{itemize}
  \item \textbf{Desindustrialización aparente:} Se observa una pérdida de peso de la industria en el VAB a precios corrientes (del 22,6\% en 1985 al 11,7\% en 2024) y en el empleo (del 19,7\% al 9,9\%). Sin embargo, medido a \textbf{precios constantes}, la caída es mucho menor (del 16,3\% al 12,2\%). Esto se explica por la mayor eficiencia industrial, la externalización de servicios (tercerización) y el impacto de la globalización.
  \item \textbf{Internacionalización y ciclo:} El sector está fuertemente internacionalizado (las exportaciones manufactureras suponen más del 20\% del PIB). Además, las manufacturas sufren fluctuaciones cíclicas mucho más violentas que el resto de la economía (caen más en recesión, suben más en expansión).
  \item \textbf{Fragmentación internacional (cadenas de valor):} Durante décadas, se deslocalizaron fases rutinarias e intensivas en trabajo hacia países de menores costes. No obstante, las tensiones geopolíticas y los cuellos de botella de la pandemia de la COVID-19 han frenado esta tendencia, propiciando una reconfiguración de las cadenas y un incipiente retorno de la producción (\textit{reshoring/nearshoring}).
\end{itemize}

\subsubsection*{3. Especialización productiva y comercial}

\begin{itemize}
  \item \textbf{Anomalía estructural:} A diferencia de Alemania (donde las manufacturas tradicionales pesan un 16\%), en España las actividades de \textbf{tecnología baja y media-baja absorben el 66,3\% del VAB industrial} (destacando la industria de alimentación y bebidas).
  \item \textbf{Déficit en alta tecnología:} La alta tecnología apenas supone el 6,4\% del VAB, lastrada por la deslocalización de la electrónica/informática hacia Asia y Europa del Este, siendo la industria farmacéutica la única excepción positiva.
  \item \textbf{Dinamismo exportador:} Irónicamente, el mayor dinamismo exportador reside en los sectores tradicionales (alimentos, cerámica, textil) y en los de tecnología media-alta (vehículos de motor y química), donde España goza de claras Ventajas Comparativas Reveladas (VCR).
  \item \textbf{Causas de esta especialización:} Reducida dimensión empresarial, abundancia relativa de trabajo poco cualificado frente a capital tecnológico, y una excesiva dependencia del capital extranjero para el desarrollo de actividades avanzadas.
\end{itemize}

\subsubsection*{4. Eficiencia productiva}

La productividad del trabajo es el termómetro de la industria.

\begin{itemize}
  \item \textbf{Evolución:} Entre 1995 y 2024, la productividad española creció a menor ritmo que la media de la Eurozona (1,4\% vs. 1,9\% anual acumulativo). Durante la burbuja (1995-2007), los Costes Laborales Unitarios (CLU) se dispararon, mermando gravemente la competitividad.
  \item \textbf{Ajuste y devaluación interna:} Durante la Gran Recesión (2008-2013), la competitividad se recuperó drásticamente mediante el cierre de empresas ineficientes, la destrucción masiva de empleo y una severa moderación salarial. Este control salarial se ha mantenido hasta la actualidad para sostener la competitividad frente a la UE.
  \item \textbf{La ``paradoja de la exportación'':} ¿Cómo puede España exportar con éxito si su eficiencia agregada es limitada? Se explica por:
    \begin{enumerate}
      \item Una fuerte heterogeneidad (coexisten pymes ineficientes enfocadas al mercado interno con gigantes multinacionales hiperproductivos).
      \item La especialización en segmentos de menor valor añadido dentro de cada industria, apoyada en salarios comparativamente más bajos que en el núcleo europeo.
      \item Diversificación hacia mercados emergentes.
      \item Una favorable relación calidad-precio.
    \end{enumerate}
\end{itemize}

\subsubsection*{5. Política industrial}

\begin{itemize}
  \item \textbf{Trayectoria:} Se ha pasado de una etapa intervencionista de ``reconversión industrial'' (años 80 y 90, con fuertes recursos para reestructurar sectores en crisis y privatizar vía SEPI) a una política industrial de corte ``horizontal'' o liberal (focalizada en I+D+i, formación e internacionalización, sin elegir ``sectores ganadores'').
  \item \textbf{Nuevo paradigma (post-COVID):} La crisis de suministros durante la pandemia demostró el gravísimo riesgo de la dependencia externa en bienes industriales esenciales (respiradores, semiconductores).
  \item \textbf{Fondos Next Generation EU:} Constituyen el eje actual de la reindustrialización, ofreciendo recursos históricos para lograr la transformación digital y la transición energética (energías renovables, vehículo eléctrico, robótica) del tejido manufacturero español.
\end{itemize}

\clearpage
% TEMA 5: El sector servicios
\section*{Tema 5: El sector servicios}

\subsection*{1. Delimitación y clasificación}

El sector terciario abarca un conjunto de actividades extraordinariamente heterogéneas (transporte, telecomunicaciones, comercio, hostelería, sanidad, educación, servicios financieros y administraciones públicas). Su delimitación se articula mediante tres clasificaciones fundamentales:

\begin{itemize}
  \item \textbf{Servicios de mercado vs. no destinados a la venta:} Los primeros son suministrados mediante el mecanismo de mercado, mientras que los segundos son provistos por el sector público de forma gratuita o a precios desvinculados de su coste de producción real.
  \item \textbf{Servicios intermedios vs. finales:} Se diferencian según si se emplean como consumos (inputs) para otras producciones o si se destinan directamente al consumidor final.
  \item \textbf{Servicios estancados vs. progresivos:} Los estancados presentan serias dificultades para sustituir mano de obra por tecnología sin mermar la calidad o cantidad; en contraste, los progresivos permiten una capitalización creciente que favorece el avance de la productividad.
\end{itemize}

Las actividades terciarias se caracterizan por su intangibilidad, la inseparabilidad espacial y temporal entre producción y consumo, su alta intensidad en mano de obra y la exigencia de contacto directo con el usuario. Estas características generan problemas analíticos severos:

\begin{itemize}
  \item En muchos servicios públicos (Sanidad, Educación), el output se mide por los inputs (costes) empleados, lo que imposibilita un cálculo preciso de la productividad.
  \item Es muy complejo cuantificar las mejoras de calidad por la fuerte subjetividad del consumidor.
  \item Existen problemas de contabilización sectorial cuando las empresas externalizan servicios que previamente realizaban de forma interna.
\end{itemize}

Finalmente, el proceso de \textbf{terciarización} (crecimiento del peso de los servicios en producción y empleo por la liberación de mano de obra de la agricultura e industria) es un fenómeno tardío en España (a partir de los años 80) en comparación con el resto de la Unión Europea.

\subsection*{2. Evolución del sector}

El sector servicios ha experimentado un crecimiento incontestable en las últimas décadas:

\begin{itemize}
  \item \textbf{Valor Añadido Bruto (VAB):} A precios corrientes, pasó de suponer el 59,5\% del total nacional en 1985 al 75,8\% en 2024. A precios constantes (base 2015), en 2024 alcanza el 76,9\%.
  \item \textbf{Empleo:} La ocupación ha pasado de representar el 57,1\% en 1985 al 78,1\% del total en 2024.
  \item \textbf{Comercio exterior:} Las exportaciones de servicios representan el 34,4\% del total de las exportaciones españolas en 2024. Hasta 1990, esta partida dependía casi en exclusiva del turismo, pero desde entonces se ha diversificado con fuerza hacia otras actividades. España se erige como una potencia mundial exportadora en servicios.
\end{itemize}

\subsection*{3. Especialización productiva y comercial}

La estructura de los servicios en España presenta dinámicas de especialización muy concretas:

\begin{itemize}
  \item Los servicios ``no destinados a la venta'' han crecido impulsados por la descentralización administrativa y el desarrollo del Estado del Bienestar desde 1975 (educación, sanidad, asistencia).
  \item En contraste con la media de la UE-27 (2023), España presenta una especialización acusada en \textbf{Hostelería} (8,9\% del VAB español frente al 3.7\% europeo), \textbf{Actividades Inmobiliarias} (15,2\% vs. 14,9\%) y \textbf{Comercio y reparación} (16,6\% vs. 15,5\%). Por el contrario, adolecemos de un menor peso en ``Información y comunicaciones'' y ``Actividades profesionales, científicas y técnicas''.
\end{itemize}

\textbf{Estructura exterior y el Turismo:}

\begin{itemize}
  \item La exportación está dominada por ``Turismo y viajes'' (46,7\% de las exportaciones de servicios en 2023, con una brutal tasa de cobertura del 323,2\%). Le siguen en importancia los ``Otros servicios empresariales'' (19,4\%).
  \item \textbf{El impacto de la pandemia:} En 2019, el turismo aportaba cerca del 12\% del PIB/empleo, recibiendo 83,5 millones de turistas extranjeros y 92,2 mil millones de euros. En 2020, el sector colapsó un 80\% (solo 18,9 millones de visitantes). Afortunadamente, en 2024 se ha completado la recuperación alcanzando la cifra récord de 93,8 millones de visitantes internacionales.
\end{itemize}

\subsection*{4. Eficiencia productiva}

La productividad de los servicios españoles es crónicamente inferior a la de la UE-27, aportando al crecimiento de la producción un 29,8\% en el periodo 2001-2023, frente al 39,8\% de la media europea.

\begin{itemize}
  \item \textbf{Ganancias de productividad (2001-2023):} Se localizan en ``Distribución comercial'' (+1,3\% gracias a hipermercados, descuento y ventas online), ``Información y comunicaciones'' (+1,5\%) y ``Finanzas y seguros'' (+1,3\%).
  \item \textbf{Descensos de productividad:} Dramático en la \textbf{Hostelería} (-0,8\% anual) por su alta intensidad laboral y reducida dimensión empresarial (impidiendo economías de escala). También en ``Sanidad y servicios sociales'' (-0,9\%) y ``Educación'' (-0,2\%).
  \item \textbf{Causas del estancamiento:} Nuestra especialización en ramas intensivas en mano de obra no cualificada (comercio, turismo), la existencia de mercados con baja competencia (que frena la innovación) y el excesivo minifundismo empresarial (reducido tamaño medio).
\end{itemize}

\subsection*{5. Política sectorial}

Históricamente, los servicios sufrieron un exceso de regulación limitante y concesiones administrativas defendidas por grupos de interés. Desde la década de 1990, impulsada por las directivas del Mercado Único Europeo, se inició un intenso proceso de \textbf{desregulación, liberalización y privatización}. Esto ha supuesto un aumento de la eficiencia, bajada de precios, mejoras en calidad y una redistribución de rentas a favor del consumidor.

La política de red (desde 1998) ha generado impactos clave por sector:

\begin{itemize}
  \item \textbf{Telecomunicaciones:} Fuerte apertura y caída de precios, aunque recientemente se ha concentrado en un oligopolio donde MasOrange, Telefónica y Vodafone controlan el 90,2\% del mercado móvil en 2024.
  \item \textbf{Transporte ferroviario:} España es pionera europea en la liberalización de pasajeros. Desde 2020-2022 operan 3 compañías distintas (Renfe/Avlo, Ouigo e Iryo).
  \item \textbf{Transporte aéreo:} La irrupción de las aerolíneas de bajo coste disparó las rutas y la competitividad.
  \item \textbf{Servicios postales:} Liderados aún por la estatal Correos, en 2023 ya había inscritos 3.119 operadores postales, forzando una mejora radical en los plazos de entrega.
  \item \textbf{Distribución comercial:} La normativa choca con diferencias legislativas autonómicas sobre horarios y grandes superficies, aunque se ha flexibilizado desde 2012.
  \item \textbf{Transporte en taxi y VTC:} La competencia (20.519 licencias VTC a 1 de enero de 2025, mayoritariamente en Madrid) ha obligado a modernizar el taxi tradicional (pagos con tarjeta, tarifas cerradas).
  \item \textbf{Farmacias:} Continúa siendo un mercado altamente protegido, con exigencias de distancias, ratios poblacionales y colegiación obligatoria, sin reformas sustanciales recientes.
\end{itemize}
 
\clearpage

% TEMA 6: Las relaciones económicas exteriores
\section*{Tema 6: Las relaciones económicas exteriores}

\subsection*{1. El sector exterior de la economía española}

La economía española ha experimentado un profundo proceso de apertura e internacionalización desde mediados del siglo XX. Este fenómeno ha sido un motor de desarrollo económico a largo plazo.

\begin{itemize}
  \item El progreso técnico se ha importado fundamentalmente a través de patentes, licencias y la entrada de inversión extranjera.
  \item La apertura comercial ha incrementado los niveles de bienestar y eficiencia general.
  \item El patrón comercial español se ha asimilado al de las economías más avanzadas. Además, se ha producido una intensa globalización productiva mediante la implantación de empresas extranjeras en España y la expansión de firmas nacionales en el exterior.
\end{itemize}

\subsection*{2. El equilibrio exterior de la Balanza de Pagos}

La Balanza de Pagos es el documento contable que registra las operaciones comerciales y financieras de los residentes de un país con el resto del mundo durante un año. Su elaboración se rige por el VI Manual de la Balanza de Pagos del FMI, adoptado por la Unión Europea en 2014. Se articula en tres grandes cuentas:

\subsubsection*{2.1. La Cuenta Corriente}

Históricamente, la cuenta corriente española ha sido deficitaria, habiendo logrado superávits únicamente en periodos muy concretos (1960-62, 1970-73, 1978-79, 1984-86, y recientemente de 2013 a 2024). Este déficit crónico ha supuesto históricamente una necesidad de financiación exterior y la consiguiente acumulación de deuda.

Sus componentes principales son:

\begin{itemize}
  \item \textbf{Balanza Comercial (Bienes):} Presenta un saldo crónicamente negativo. Esta posición estructuralmente importadora obedece principalmente a las compras masivas de bienes energéticos.
  \item \textbf{Rentas Primarias:} Muestra un saldo habitualmente deficitario debido a que España es un país receptor neto de capital foráneo, obligando al pago de dividendos e intereses al exterior. Tras unos años de equilibrio relativo (2015-2022), el déficit reapareció en 2023 y 2024 a causa de las subidas de los tipos de interés.
  \item \textbf{Rentas Secundarias:} Registró superávit hasta 1985 gracias a las remesas de emigrantes españoles. Posteriormente, se estabilizó en un déficit de en torno al 1\% del PIB.
  \item \textbf{Servicios:} Constituye la gran partida compensatoria, arrojando un superávit muy estable (promedio del 3\% del PIB entre 1985 y 2024). El turismo y los servicios empresariales y tecnológicos han catapultado sus ingresos totales hasta representar el 12\% del PIB en 2024.
\end{itemize}

\subsubsection*{2.2. La Cuenta de Capital}

Registra las transferencias internacionales de capital y la compraventa de activos intangibles.

\begin{itemize}
  \item El grueso de los ingresos proviene de los fondos comunitarios estructurales de la UE (como el FEADER y el FEDER).
  \item El saldo ha sido positivo (~1\% del PIB) desde 1993, experimentando una caída en 2005 por la ampliación de la UE hacia el Este, y un notable repunte desde 2022 propulsado por los fondos europeos \textit{Next Generation EU}.
\end{itemize}

\subsubsection*{2.3. La Cuenta Financiera}

Refleja la Variación Neta de Activos (VNA) y la Variación Neta de Pasivos (VNP) frente al exterior.

\begin{itemize}
  \item Tradicionalmente, los déficits por cuenta corriente forzaban una entrada neta de capital extranjero (VNP > VNA).
  \item Con la crisis de deuda soberana (2011-2012), se produjo una salida neta de fondos de España (caída abrupta de inversión externa), la cual tuvo que ser sostenida mediante los pasivos del Banco de España frente al Eurosistema.
  \item Afortunadamente, a partir de 2013 la tendencia se invirtió, logrando la economía española capacidad de financiación neta frente al mundo (VNA > VNP).
\end{itemize}

\subsubsection*{2.4. El Ajuste del Déficit Exterior}

El déficit corriente ha sido un cuello de botella sistémico.

\begin{itemize}
  \item Antes del euro, los gobiernos ajustaban los fuertes déficits comerciales aplicando políticas de enfriamiento de la demanda y, en última instancia, recurriendo a devaluaciones competitivas de la peseta.
  \item Al ingresar en la Unión Económica y Monetaria, se perdió la herramienta cambiaria. Durante la burbuja (1995-2007), la abundancia de capital extranjero barato alimentó la demanda interna hasta generar un déficit corriente insostenible del 10\% del PIB en 2007.
  \item Desde entonces se operó un severo ajuste, transitando hacia superávits constantes que han escalado al 4\% del PIB en 2023 y 2024.
\end{itemize}

\subsection*{3. El comercio exterior de España}

El crecimiento económico de las últimas décadas se asienta en el despegue comercial.

\begin{itemize}
  \item \textbf{Apertura y cobertura:} Las exportaciones e importaciones de bienes han pasado de representar conjuntamente el 32\% del PIB en 1980 a suponer un grado de apertura del 73\% del PIB en 2024 (incluyendo los servicios). Nuestra tasa de cobertura comercial (exportaciones sobre importaciones) superó el 90\% para bienes desde 2013, y excede holgadamente el 100\% si incorporamos los servicios.
  \item \textbf{Comportamiento elástico:} Nuestras exportaciones dependen de la renta de los países clientes (elasticidad cercana a 1) y de nuestros precios relativos internacionales. Las importaciones, sin embargo, reaccionan violentamente a la renta nacional (elasticidad = 2). Por ello, cuando España crece más rápido que la media de la OCDE, nuestro saldo exterior tiende inexorablemente a deteriorarse.
  \item \textbf{Impacto pre y post 2008:} El enorme déficit previo a 2008 se debió a un fuerte crecimiento diferencial respecto a la OCDE, una inflación más alta y la escalada del crudo. El drástico ajuste post-crisis se logró mediante la contención de las importaciones por el desplome interno y una fuerte ``devaluación interna'' (reducción de salarios) que impulsó la competitividad exterior.
  \item \textbf{Orientación geográfica y sectorial:} La Unión Europea es nuestro mercado vital, absorbiendo el 62,6\% de nuestras exportaciones y originando el 49,3\% de nuestras importaciones en 2023. Sufrimos déficit recurrente frente a China y los países de la OPEP. Sectorialmente, nuestro comercio domina en manufacturas de tecnología media-alta. Poseemos ventajas relativas en los bienes de consumo y automoción, mientras arrastramos fuertes desventajas en la alta tecnología.
  \item \textbf{Comercio intraindustrial:} Destaca el intercambio bidireccional de bienes similares con los países de la UE, que ha escalado hasta el 70\% impulsado por la semejanza de renta y las economías de escala.
\end{itemize}

\subsection*{4. Las Inversiones Extranjeras Directas (IED)}

Se define la IED como aquella inversión donde se adquiere al menos el 10\% del capital social para ejercer el control y la gestión continua de la firma, a diferencia de la inversión de cartera, que persigue una rentabilidad especulativa a corto plazo.

\begin{itemize}
  \item \textbf{España como país receptor:} Tras integrarse en la UE, la IED recibida se triplicó rápidamente (1986-1995). Pese a las crisis, hemos consolidado el \textit{stock} extranjero. En 2023, España se ubica como el 13º receptor mundial (1,8\% del total global), destacando netamente por encima de nuestro peso en el PIB mundial. Los sectores más atractivos han sido los servicios (49,2\% del total, especialmente inmobiliaria y finanzas), las manufacturas tradicionales (24,1\%) y el sector energético (17,2\%), este último disparado por los ingentes proyectos \textit{greenfield} en energías renovables. El capital proviene mayormente de Estados Unidos, Reino Unido, Francia y Alemania.
  \item \textbf{España como país emisor:} Las firmas españolas protagonizaron un despegue internacional épico. En 1985 el \textit{stock} de IED emitida apenas sumaba el 2,5\% de nuestro PIB. Al calor del boom expansivo y el euro, ese esfuerzo saltó al 39,5\% del PIB en 2007. Hoy (2023), España se afianza como el 17º inversor global, destinando un colosal 65,8\% de su IED a servicios (liderados por el sector bancario y asegurador), canalizados predominantemente hacia la Eurozona, América del Norte y Latinoamérica (con México y Brasil a la cabeza).
\end{itemize}


\clearpage

% TEMA 7: Mercado de trabajo
\section*{Tema 7: Mercado de trabajo}

\bigskip
\subsection*{1. Caracterización del mercado de trabajo en España}

El mercado laboral español, desde 1980 hasta la actualidad, presenta rasgos muy singulares en su evolución y estructura.

\begin{itemize}
  \item \textbf{Evolución poblacional y de empleo:} Se observa un aumento histórico de la población activa impulsado por el incremento de la población en edad de trabajar, la inserción laboral femenina y la inmigración (especialmente en la primera década de los 2000).

  \item \textbf{Volatilidad cíclica:} Existe un elevado ritmo de creación de empleo en fases expansivas (como entre mediados de los 90 y 2007, o a partir de 2014) que contrasta drásticamente con la fuerte destrucción de puestos en etapas recesivas (crisis de los 70, 90 y la Gran Recesión de 2008).

  \item \textbf{Comparativa con la Unión Europea (UE-27):}
    \begin{itemize}
      \item \textbf{Tasa de desempleo:} Es significativamente más elevada en España, aunque se redujo fuertemente hasta 2007 y tras 2014.
      \item \textbf{Ocupación y Actividad:} La tasa de ocupación es menor que en la UE debido al impacto de crisis pasadas. La tasa de actividad se sitúa levemente por debajo de la media europea, lastrada recientemente por la menor participación masculina.
      \item \textbf{Calidad del empleo:} España sufre de una alta tasa de temporalidad que supera la media europea en periodos de recuperación. El empleo a tiempo parcial es escaso y frecuentemente involuntario.
      \item \textbf{Formación:} Destaca un porcentaje superior a la media europea de empleados en formación continua.
    \end{itemize}

  \item \textbf{Perfil del desempleado:} La incidencia del paro es severa entre los jóvenes (duplicando la tasa media), mayor en las mujeres (brecha que volvió a crecer tras 2014), inmigrantes y personas con bajo nivel educativo. Geográficamente, existen fuertes disparidades entre el norte (menor paro) y el sur/Canarias (mayor paro) debido a la escasa movilidad y distintas estructuras productivas.
\end{itemize}

\subsection*{2. Marco institucional del mercado de trabajo y reformas laborales}

Las instituciones y normativas definen el funcionamiento de nuestro mercado laboral. Históricamente, desde el Estatuto de los Trabajadores (1980), se implantó un marco similar al europeo pero con fuertes rigideces.

Existen cuatro grandes bloques institucionales:

\begin{enumerate}
  \item \textbf{Fijación de salarios:} Depende de la negociación colectiva. Pese a la baja afiliación sindical, la cobertura de los convenios es altísima. Históricamente, predominaba un nivel de negociación intermedio (sectorial), lo que generaba menor moderación salarial y, por ende, mayor paro.

  \item \textbf{Costes de despido (flexibilidad externa):} Históricamente altos en España (indemnizaciones que doblaban la media europea), lo que provocó que las empresas recurrieran masivamente a la contratación temporal para ajustar plantillas, generando un mercado dual.

  \item \textbf{Prestaciones por desempleo:} Cuanto mayor es la tasa de sustitución y la duración, menor es la intensidad de búsqueda de empleo. Las reformas de 1992 y 2012 redujeron su generosidad para alinearla al promedio europeo.

  \item \textbf{Políticas Activas de Empleo:} Destinadas a formar y asesorar al parado. España gasta menos en ellas en proporción a su paro respecto a la UE, priorizando tradicionalmente subvencionar la contratación pasiva frente a la recualificación real.
\end{enumerate}

\clearpage
% TEMA 9: El sector público de la economía española
\section*{Tema 9: El sector público de la economía española}

\subsection*{1. El papel del Estado en las economías actuales}

El sector público es una pieza fundamental en las economías desarrolladas contemporáneas. La crisis de 1929 evidenció la incapacidad de los mercados para recuperarse de forma autónoma. Posteriormente, la \textit{Teoría General} de Keynes (1936) subrayó la necesidad de utilizar la política monetaria y fiscal, lo que disparó el peso del sector público sobre el PIB en Europa occidental tras la Segunda Guerra Mundial.

Así se consolidaron las economías mixtas, donde el mercado es el motor principal, pero el sector público asume un papel relevante corrigiendo y complementando sus deficiencias mediante instituciones y normas. Históricamente, la corriente liberal de los años ochenta frenó este crecimiento en Europa, aunque en España el desarrollo público fue más tardío y acelerado. Recientemente, la crisis de 2008 exigió duros recortes de gasto y subidas de impuestos ante el incremento del déficit, mientras que la pandemia de COVID-19 devolvió al Estado un papel protagonista mediante la protección social y el apoyo empresarial.

El Estado asume funciones insoslayables en la economía:

\begin{itemize}
  \item Mantener un marco institucional que fije las ``reglas del juego'' (propiedad, leyes antimonopolio) y corrija fallos de mercado.
  \item Intervenir mediante regulaciones concretas de precios o cantidades.
  \item Actuar como agente económico puro, financiándose o realizando actividad empresarial en los mercados.
  \item Utilizar la Hacienda Pública (impuestos y gastos) para redistribuir la renta y suavizar los ciclos económicos.
  \item Promover una asignación eficiente de los recursos acorde a las necesidades sociales.
\end{itemize}

\subsection*{2. El sector público en España: un Estado descentralizado}

La delimitación del sector público admite dos perspectivas: una amplia, que incluye a las Administraciones Públicas (AAPP) y a las Empresas Públicas, y una estricta, la más habitual, que engloba exclusivamente a las Administraciones Públicas.

Las AAPP operan bajo el criterio de autoridad, ofrecen servicios no destinados a la venta y se financian coercitivamente mediante tributos. Por su parte, las empresas públicas actúan en el mercado persiguiendo fines estratégicos y se financian vendiendo bienes o servicios.

El organigrama de las AAPP en España se estructura de la siguiente manera:

\begin{itemize}
  \item \textbf{Unión Europea:} Establece regulaciones que impactan fuertemente en la economía y la hacienda española.
  \item \textbf{Administración Central:} Compuesta por el Estado y los Organismos Autónomos (como el INE o el SEPE).
  \item \textbf{Administraciones Territoriales:} Integradas por las Comunidades Autónomas (CC.AA.) y las Corporaciones Locales.
  \item \textbf{Seguridad Social:} Administra los gastos de protección social (pensiones, sanidad) con una financiación mixta de cotizaciones y transferencias estatales.
\end{itemize}

España destaca por ser uno de los países más descentralizados de Occidente. Este proceso se originó con la Constitución de 1978, que propició un masivo trasvase de competencias hacia las CC.AA. y la Seguridad Social. En 2024, el 44\% del gasto correspondía a las Administraciones Territoriales y el 56\% al nivel Central (inflado por el peso prestacional de la Seguridad Social, no por recentralización de funciones). Asimismo, el empleo público se concentra actualmente en las CC.AA. con un 62\%, frente al 20\% local y el 18\% central. Las empresas públicas estatales hoy son residuales (1,4\% del PIB), aunque las autonómicas han ganado terreno en infraestructuras y servicios.

\subsection*{3. El marco regulador}

La tradición económica española fue marcadamente intervencionista, agudizándose durante el franquismo. En 1975, España padecía una hacienda débil y una regulación asfixiante sobre los factores productivos y el comercio.

Esta dinámica cambió a través de dos ejes:

\begin{enumerate}
  \item \textbf{Reforma fiscal:} Impulsada por los Pactos de la Moncloa y la Constitución, dotó a la hacienda de suficiencia y equidad.
  \item \textbf{Desregulación:} Acelerada por la integración europea y el Mercado Único. Hoy, dos tercios de las normativas económicas provienen de la UE.
\end{enumerate}

El paradigma viró desde las regulaciones de estructura (barreras de entrada) hacia regulaciones de conducta (supervisión del comportamiento), creando entidades supervisoras como la actual Comisión Nacional de los Mercados y la Competencia (CNMC). Sin embargo, persisten graves problemas: el marco es hipercomplejo, las normas son de escasa calidad, existe falta de coordinación entre los niveles de gobierno y todo ello encarece los costes empresariales, reduciendo la productividad y generando incertidumbre.

\subsection*{4. La Hacienda de las Administraciones Públicas}

La Hacienda Pública ejerce su impacto a través de sus ingresos y gastos. Actúa mediante \textbf{variaciones discrecionales} (decisiones activas del gobierno, como modificar tipos impositivos o reasignar partidas) y \textbf{estabilizadores automáticos} (mecanismos que se ajustan solos a la coyuntura, como la recaudación de impuestos progresivos o el aumento de prestaciones por desempleo en las crisis).

\subsubsection*{4.1. Gasto Público y Estado del Bienestar}

La dimensión del sector público se mide mediante el cociente de Gasto Público sobre el PIB. En España, este indicador pasó del 23,5\% en 1975 al 44,9\% en 2024. Esto sitúa a nuestra economía en un modelo intermedio, por encima de Estados Unidos (40,1\%) pero inferior a la media de la Eurozona (49,4\%). Este crecimiento responde a la descentralización, la consolidación del Estado del Bienestar y los estabilizadores automáticos.

La trayectoria temporal del gasto muestra repuntes en la crisis de los 90, una consolidación de 10 puntos de caída entre 1995 y 2007 (previa a la Unión Monetaria), y fuertes aumentos por las crisis de 2008 (elevándolo 7 puntos) y de la pandemia en 2020 (alcanzando el 51,4\% del PIB), para reconducirse al 45\% entre 2023 y 2024.

\textbf{Destino final de los recursos (2024):}

\begin{itemize}
  \item \textbf{Estado del Bienestar (67,2\%):}
    \begin{itemize}
      \item Las pensiones absorben el 31,1\% debido a la dinámica demográfica y bases reguladoras más altas indexadas al IPC. Las reformas de 2011/2013 retrasaron la edad a 67 años y ampliaron el cómputo. Las recientes de 2021/2023 vincularon el alza al IPC e introdujeron el mecanismo de equidad intergeneracional.
      \item La salud representa el 14,5\%, tensionada por el envejecimiento, las nuevas tecnologías y las secuelas del COVID-19.
      \item La educación acapara un 9,3\%, estancada por la caída de la natalidad.
      \item El desempleo supone el 3,4\%.
    \end{itemize}
  \item \textbf{Servicios Públicos (18,5\%):} Incluye los intereses de la deuda (2,5\% del PIB), llamativamente altos a pesar de las políticas laxas del BCE previas.
  \item \textbf{Asuntos Económicos (11,0\%):} En aumento reciente por los fondos \textit{Next Generation EU} orientados a la transición digital y ecológica.
\end{itemize}

Comparativamente con la UE-27, España destina menos recursos a educación, sanidad, defensa y protección social familiar, invirtiendo comparativamente más en pensiones y desempleo. En su distribución competencial, el Estado financia los bienes públicos puros; la Seguridad Social, las transferencias monetarias de protección; las CC.AA., los servicios sanitarios y educativos; y los entes locales, los servicios comunitarios de proximidad.

\subsubsection*{4.2. Ingresos Públicos}

El sistema tributario contemporáneo se gestó con la reforma de 1977 y maduró en 1986 con la implantación del IRPF, el IVA y Patrimonio.

Estructura de ingresos (2023):

\begin{itemize}
  \item \textbf{Cotizaciones sociales (31,4\%):} Eje vertebrador del pago de pensiones. Las reformas de 2021 y 2023 han ajustado sus tipos, topes máximos y la cuota de autónomos por ingresos reales.
  \item \textbf{Impuestos directos (30,0\%):} Muy procíclicos. Destacan el IRPF y Sociedades. Como apunte, el impuesto de patrimonio apenas recauda un 0,1\% del PIB.
  \item \textbf{Impuestos indirectos (26,3\%):} El IVA aporta más de la mitad de esta partida, seguido por los impuestos especiales sobre hidrocarburos, tabaco y alcohol.
\end{itemize}

La \textbf{Presión Fiscal} (Impuestos y Cotizaciones sobre el PIB) ha escalado desde el 35,7\% en 1985 hasta el 41,9\% en 2024. Pese a este incremento (especialmente acusado entre 2021 y 2023 frente al estancamiento de la UE-27), España mantiene una presión fiscal inferior a la media comunitaria, apoyándose comparativamente más en imposición directa y cotizaciones, y menos en fiscalidad indirecta.

\subsection*{5. Saldo Presupuestario y Deuda}

La historia del saldo público dibuja tres épocas claramente diferenciadas:

\begin{enumerate}
  \item \textbf{1975-1995:} Expansión crónica del déficit por la construcción del Estado del Bienestar y las deficiencias recaudatorias en las crisis.
  \item \textbf{1996-2007:} Etapa de saneamiento lograda gracias a mayor presión fiscal, moderación salarial pública y bajos intereses, culminando en los superávits inéditos del periodo 2005-2007.
  \item \textbf{2008-Actualidad:} Deterioro monumental. El déficit tocó fondo en 2009 (-11\% del PIB). Hubo un repunte anómalo en 2012 por el rescate bancario (-3,7\%), lográndose reducir hasta el 2,6\% en 2018. Posteriormente, la pandemia hundió el saldo al -10\% en 2020, que se ha corregido paulatinamente hasta el -3,2\% previsto para 2024.
\end{enumerate}

Resulta clave distinguir analíticamente entre el déficit cíclico y el \textbf{déficit estructural} (aquel que existiría si la economía estuviese en su PIB potencial). España padece un severo déficit estructural crónico producto de políticas discrecionales y no puramente del ciclo recesivo. Esta losa se infló en 2008 y frenó su corrección en 2015 debido a rebajas de impuestos centrales y aumentos de gastos autonómicos.

Finalmente, este descuadre recurrente se financia emitiendo \textbf{Deuda Pública}. Medida sobre el PIB, la deuda descendió virtuosa y sostenidamente desde 1995 (61,7\%) hasta 2007 (35,6\%). Sin embargo, la concatenación de la Gran Recesión y la pandemia fulminó ese progreso, catapultándola por encima del 100,4\% en 2014, el 119,3\% en 2020 y moderándose ligeramente al 103,5\% en 2024. Además del volumen de pasivo, las crisis detonaron tensiones severas en los costes de financiación, materializadas dramáticamente en la prima de riesgo que en 2012 llegó a exigir un diferencial del 6\% frente al bono alemán.
 
\subsubsection*{Reformas Laborales Clave (Exigidas para el estudio)}

\begin{itemize}
  \item \textbf{Reforma de 2012:} Buscó atajar el paro estructural reduciendo rigideces. Descentralizó la negociación dando prioridad al convenio de empresa y limitando la ultraactividad. Redujo el coste del despido improcedente (de 45 a 33 días) y facilitó los despidos objetivos. Permitió que los salarios reaccionaran mejor al desempleo.

  \item \textbf{Reforma de 2021:} Su objetivo primordial fue reducir la temporalidad y precariedad. Estableció el contrato indefinido por defecto y eliminó el de ``obra o servicio''. Limitó la temporalidad a circunstancias de producción o sustitución. Potenció la figura del ``fijo discontinuo''. Revirtió partes de 2012, recuperando la ultraactividad y la prioridad del convenio sectorial en materia salarial. Creó los nuevos ERTEs y el Mecanismo RED para crisis cíclicas o sectoriales.
\end{itemize}

\subsection*{3. El desempleo y sus causas}

El desempleo se divide en dos componentes que en España son singularmente altos:

\begin{itemize}
  \item \textbf{Desempleo Estructural (NAIRU):} Tasa de desempleo compatible con una inflación estable. Surgió en los años 70 por shocks de oferta (petróleo, caída de productividad, alzas salariales) combinados con instituciones rígidas. Sus causas incluyen: rigidez salarial (los salarios no bajaban aunque hubiera paro), baja intensidad de búsqueda de empleo y desajustes geográficos y formativos. La reforma de 2012 logró reducir esta NAIRU en casi 4 puntos porcentuales.

  \item \textbf{Desempleo Cíclico y Temporalidad:} Producto de variaciones en la demanda agregada. España tradicionalmente ajustó sus crisis destruyendo empleo temporal (ajuste por cantidades) en lugar de ajustar salarios o jornadas (flexibilidad interna).

  \item \textbf{Impacto de la COVID-19:} Se evidenció un ``desacoplamiento''. Aunque el PIB y las horas efectivas cayeron en torno al 11\%, el empleo oficial solo bajó un 2,9\% gracias al ``escudo'' de los ERTEs, que protegieron a 3,6 millones de trabajadores, evitando la destrucción masiva.

  \item \textbf{Evaluación y retos actuales:} La reforma de 2021 redujo estadísticamente la temporalidad al 15,5\%, pero generó una ``ilusión contractual'': persiste la rotación bajo nuevos formatos y los despidos en periodos de prueba. Los retos pendientes son el alto subempleo involuntario, reformar las políticas activas hacia la recualificación (ALMPs) y solucionar la ``Paradoja de la escasez'' (puestos vacantes en sectores intensivos pese al alto paro).
\end{itemize}

\subsection*{4. Indicadores básicos del mercado de trabajo}

El análisis del mercado laboral exige precisión cuantitativa. La población de 16 y más años se divide en Inactivos (estudiantes, jubilados) y Activos, y estos últimos en Ocupados (asalariados o cuenta propia) y Desempleados.

Las tasas fundamentales de análisis se formulan de la siguiente manera:

\[
\text{Tasa de Actividad} = \frac{\text{Poblacion activa}}{\text{Poblacion de 16 y mas anos}} \times 100
\]

\[
\text{Tasa de Empleo} = \frac{\text{Poblacion ocupada}}{\text{Poblacion de 16 y mas anos}} \times 100
\]

\[
\text{Tasa de Paro} = \frac{\text{Poblacion parada}}{\text{Poblacion activa}} \times 100
\]

\[
\text{Tasa de Asalarizacion} = \frac{N^{o} \text{ de asalariados}}{\text{Poblacion ocupada}} \times 100
\]

\[
\text{Tasa de Temporalidad} = \frac{N^{o} \text{ asalariados con contrato temporal}}{N^{o} \text{ total de asalariados}} \times 100
\]

\[
\text{Tasa de Trabajo a tiempo parcial} = \frac{N^{o} \text{ ocupados a tiempo parcial}}{\text{Poblacion ocupada}} \times 100
\]

\textbf{Productividad y Costes Laborales:}
La Productividad Aparente del Trabajo se define como:

\[
\text{PAT} = \frac{\text{PIB}_{\text{nominal}}}{\text{Empleo}}
\]

Para medir los costes laborales se emplean los Costes Laborales Unitarios (CLU):

\[
\text{CLU}_{\text{nominal}} = \frac{\frac{\text{RA}}{\text{Asalariados}}}{\frac{\text{PIB}_{\text{nominal}}}{\text{Empleo}}}
\]

\textit{(Donde RA es Remuneración de Asalariados)}

\[
\text{CLU}_{\text{real}} = \frac{\text{CLU}_{\text{nominal}}}{\text{Deflactor Implicito del PIB}}
\]

Le sugiero repasar detenidamente las relaciones institucionales y los efectos macroeconómicos de las dos últimas grandes reformas, tal como se especifica en las pautas del documento. Que tenga un buen estudio.
Le sugiero repasar detenidamente las relaciones institucionales y los efectos macroeconómicos de las dos últimas grandes reformas, tal como se especifica en las pautas del documento.

\clearpage
% TEMA 8: Ámbito monetario y financiero
\section*{Tema 8: El ámbito monetario y financiero de la economía española}

 
\subsection*{1. Definición y funciones del Sistema Financiero}

El sistema financiero se define como el conjunto de instituciones, intermediarios y mercados encargados de canalizar el ahorro desde las unidades económicas con exceso de fondos hacia aquellas que carecen de financiación. Esta labor de intermediación es vital y cumple tres funciones básicas:

\begin{itemize}
  \item Transferencia de fondos entre los agentes económicos.
  \item Transferencia de riesgos, dado que los ingresos futuros del prestatario son inciertos.
  \item Transformación de plazos, compatibilizando las preferencias temporales, generalmente distintas, de ahorradores e inversores.
  \item Su papel es fundamental para el crecimiento económico, ya que provee recursos para la inversión productiva y garantiza financiación a las actividades que generan un mayor rendimiento.
\end{itemize}

\subsection*{2. Instituciones, Instrumentos y Mercados Financieros}

El marco operativo se divide en varios componentes esenciales.

\textbf{Instituciones:}
\begin{itemize}
  \item \textbf{Instituciones Monetarias:} Se caracterizan porque parte de sus pasivos funcionan como dinero o medio de pago (billetes, monedas, depósitos). Incluyen al Banco de España, los fondos del mercado monetario, las entidades de dinero electrónico, la banca privada, las cajas de ahorro y las cooperativas de crédito.
  \item \textbf{Instituciones No Monetarias:} Su actividad es estrictamente de mediación financiera. Engloban a empresas de seguros, fondos de pensiones y otros intermediarios.
\end{itemize}

\textbf{Activos Financieros:}
Son los títulos creados por las instituciones, representando un pasivo para el emisor y un activo para el poseedor. Tienen tres características fundamentales:
\begin{itemize}
  \item \textbf{Liquidez:} La facilidad de conversión en dinero sin sufrir pérdidas.
  \item \textbf{Rentabilidad:} Los rendimientos o intereses que generan.
  \item \textbf{Riesgo:} La probabilidad de que el prestatario cumpla con sus obligaciones.
\end{itemize}

\textbf{Mercados Financieros:}
Se clasifican según varios criterios:
\begin{itemize}
  \item \textbf{Fase de negociación:} Mercados primarios (creación de activos) y mercados secundarios (negociación de activos ya creados).
  \item \textbf{Tipo de títulos:} Mercados monetarios (corto plazo, alta liquidez, riesgo reducido; base de los tipos de interés) y mercados de capitales (medio y largo plazo; subdivididos en renta fija y renta variable).
  \item \textbf{Intervención de autoridades:} Libres (determinados por oferta/demanda) o regulados (con limitaciones en tipos o volúmenes obligatorios).
  \item \textbf{Grado de organización:} Organizados (con agentes oficiales en un centro único, como la Bolsa) o no organizados (intercambio directo sin lugar central).
\end{itemize}

\textbf{Modelos de Financiación:}
Se contrasta el \textit{modelo anglosajón}, con un predominio de la financiación directa en mercados, y el \textit{modelo continental} (aplicable a España), dominado por la financiación intermediada mediante un sistema de ``banca universal'' estrechamente vinculada a las empresas.

\subsection*{3. Rasgos Básicos del Sistema Financiero Español}

El sistema financiero nacional presenta particularidades estructurales y evolutivas significativas:

\begin{itemize}
  \item \textbf{Evolución:} Ha habido un desarrollo de intermediarios no bancarios (fondos, aseguradoras), aunque su peso sigue siendo menor que en países anglosajones. Como evidencia, el efectivo y depósitos bancarios suponían el 63\% de los activos de las familias en 1985, cayendo al 35\% en 2024.
  \item \textbf{Sistema Bancario:} Destaca por su gran tamaño (activos en torno al 2,7\%-2,8\% del PIB) y la reducida cuota de la banca extranjera. Está compuesto por bancos (sociedades anónimas que reparten dividendos), cajas de ahorro (entidades fundacionales sin ánimo de lucro; pasaron de 47 en 2005 a solo 2 en la actualidad) y cooperativas de crédito (cajas rurales).
  \item \textbf{Desregulación (Fines del S. XX):} La eliminación de trabas, los avances en telecomunicaciones y el desarrollo del mercado forzaron nuevas estrategias bancarias. Los bancos apostaron por ganar tamaño y expandirse en Europa, reduciendo red nacional, mientras que las cajas se expandieron nacionalmente, logrando superar a los bancos en número de oficinas y empleados a partir de 1995 (hasta 2008).
\end{itemize}

\subsection*{4. Crisis Bancaria y Reestructuración}

La crisis de 2008 transformó radicalmente el panorama bancario.

\begin{itemize}
  \item \textbf{Impacto:} Inicialmente el impacto directo fue bajo (no se permitían vehículos opacos fuera de balance), pero el impacto indirecto fue devastador debido a la alta dependencia exterior (cierre de mercados mayoristas), el desplome inmobiliario (que impactó duramente a las cajas), la caída de la actividad económica y el aumento de la prima de riesgo (deuda soberana).
  \item \textbf{Medidas iniciales:} Incluyeron inyecciones de liquidez (BCE y FAAF nacional), fomento de fusiones (FROB), la ``bancarización'' forzosa de las cajas mediante el R.D.L. 11/2010 (obligadas a ceder su negocio a un banco o convertirse en fundaciones), y exigencias superiores de provisiones y capital.
  \item \textbf{Rescate Europeo (2012):} Tras agotar el FROB, España solicitó rescate al Eurogrupo. A cambio de 100.000 millones de euros, se impusieron estrictas condiciones: recapitalización rigurosa, imposición de pérdidas a tenedores de preferentes, reducción de tamaño/plantilla, cese del negocio promotor y la creación de un ``banco malo'' (la SAREB) para aislar los activos tóxicos.
  \item \textbf{Resultado actual:} El sistema quedó fuertemente concentrado en unos pocos grandes bancos. Sufrió un severo ajuste, cerrándose unas 28.000 sucursales y eliminándose 100.000 empleos entre 2008 y 2024. Hoy, la ratio de capital de alta calidad asciende al 13\%, el doble que antes de la crisis.
\end{itemize}

\subsection*{5. Evolución Reciente}

A partir de 2014, el sector volvió a presentar beneficios gracias a la recuperación económica, reducción de morosidad y aumento de ingresos por comisiones.

\begin{itemize}
  \item \textbf{Impacto de la COVID-19:} Obligó a grandes dotaciones, pero la banca partía de una sólida posición de liquidez y capital. La intervención del BCE (inyección de 1,3 billones y financiación barata) y del Estado español (avales ICO y ERTEs) contuvo la crisis y previno el default masivo.
  \item \textbf{Escenario Post-Pandemia:} Las fuertes subidas de tipos del BCE desde 2022 elevaron significativamente los márgenes de interés bancarios. Además, el uso masivo de canales digitales aceleró el cierre de sucursales, mientras la banca adapta sus estructuras para competir contra \textit{Fintechs} y el potencial ingreso de \textit{Bigtechs}.
\end{itemize}

\subsection*{ANEXO: La Unión Bancaria Europea}

Es un proyecto estructural para reforzar la estabilidad financiera regional, de adscripción obligatoria para los países de la Eurozona. Su objetivo es asegurar entidades sólidas, evitar que los rescates recaigan en los contribuyentes y lograr la integración financiera. Se asienta sobre el ``Código Normativo Único'' de la Autoridad Bancaria Europea y tres pilares institucionales:

\begin{enumerate}
  \item \textbf{Mecanismo Único de Supervisión (MUS):} Coordinado por el BCE en conjunto con supervisores nacionales, fiscaliza las entidades más grandes.
  \item \textbf{Mecanismo Único de Resolución (MUR):} Destinado a quebrar y liquidar bancos inviables ordenadamente sin coste para el contribuyente, apoyado por el Fondo Único de Resolución (80.000 millones en 2023).
  \item \textbf{Sistema Europeo de Seguro de Depósitos (SESD):} Instrumento, aún en fase de discusión, para armonizar los fondos de garantía a nivel europeo (cubriendo hasta 100.000 euros por depositante).
\end{enumerate}